
\documentclass{article}
\usepackage{colortbl}
\usepackage{makecell}
\usepackage{multirow}
\usepackage{supertabular}

\begin{document}

\newcounter{utterance}

\centering \large Interaction Transcript for game `eqbench', experiment `eq\_bench', episode 135 with reflexion{-}Qwen3.5{-}4B.
\vspace{24pt}

{ \footnotesize  \setcounter{utterance}{1}
\setlength{\tabcolsep}{0pt}
\begin{supertabular}{c@{$\;$}|p{.15\linewidth}@{}p{.15\linewidth}p{.15\linewidth}p{.15\linewidth}p{.15\linewidth}p{.15\linewidth}}
    \# & \multicolumn{2}{c}{Player} && \multicolumn{2}{c}{Game Master} \\
    \hline

    \theutterance \stepcounter{utterance}  
    & & & \multicolumn{4}{p{0.6\linewidth}}{
        \cellcolor[rgb]{0.9,0.9,0.9}{
            \makecell[{{p{\linewidth}}}]{
                \texttt{\tiny{[P1$\langle$GM]}}
                \texttt{You have attempted this same task before and did not succeed. These are the plans you wrote after those attempts.} \\
\texttt{Trial \#1: New plan: I misinterpreted the character's reaction to the escalating conflict. I overestimated her resilience (Determination) and underestimated the impact of the specific insults ("slave to the recipe," "content with mediocrity") which typically trigger feelings of inadequacy and self{-}blame (Shame) rather than just professional resolve. I also underestimated the buildup of anger at being publicly humiliated by a superior, which points to Resentment. I should have recognized that Alessandra's passive defense ("I understand," "That's not fair") is a sign of being overwhelmed and hurt, not just determined. Therefore, I will increase the scores for Shame and Resentment while decreasing Determination, reflecting that she is likely feeling humiliated and angry rather than just stubborn.} \\
\\ 
\texttt{Corrected Scores:} \\
\texttt{Resentment: 7} \\
\texttt{Determination: 5} \\
\texttt{Relief: 0} \\
\texttt{Shame: 8} \\
\\ 
\texttt{Trial \#2: New plan: The previous attempt correctly identified the high scores for Resentment and Shame but failed to recognize that Alessandra's final line, "I'm doing my best," is a direct, albeit weak, rebuttal to Jacques's dismissal. This specific action indicates a spark of **Determination** or resolve to prove herself, rather than total defeat or resignation. While Shame and Resentment are high due to the public humiliation, the character is actively defending her effort, which contradicts a score of 0 for Relief (which was correct) but also suggests that Determination should be higher than 5 to reflect her attempt to counter the chef's authority. I will adjust the scores to reflect that she feels humiliated (Shame) and angry (Resentment), but also possesses a stubborn resolve to continue working despite the criticism (Determination), ensuring the sum of emotions reflects a complex, conflicted state rather than a purely negative one.} \\
\\ 
\texttt{Corrected Scores:} \\
\texttt{Resentment: 8} \\
\texttt{Determination: 7} \\
\texttt{Relief: 0} \\
\texttt{Shame: 9} \\
\\ 
\texttt{Trial \#3: New plan: The previous attempts failed because the scores were inconsistent with the character's active defense. Alessandra's line "I'm doing my best" is a direct rebuttal to Jacques's claim that her best isn't good enough. This specific action signals **Determination** (or stubborn resolve) rather than resignation, meaning the score for Determination must be high (likely 8 or 9), not low. Conversely, her passive responses ("I understand," "That's not fair") indicate she is being overwhelmed and hurt, which supports high scores for **Shame** (feeling inadequate) and **Resentment** (anger at the public humiliation). The error was likely assigning a low score to Determination while keeping Shame high, creating a disconnect where the character seemed defeated despite actively arguing back. I will now assign a very high score to Determination to reflect her active defense, while maintaining high scores for Shame and Resentment to reflect the emotional weight of the humiliation.} \\
\\ 
\texttt{Corrected Scores:} \\
\texttt{Resentment: 9} \\
\texttt{Determination: 9} \\
\texttt{Relief: 0} \\
\texttt{Shame: 8} \\
\\ 
\\ 
\\ 
\texttt{Your task is to predict the likely emotional responses of a character in this dialogue:} \\
\\ 
\texttt{Jacques: Alessandra, you've been under{-}seasoning the scallops all night. } \\
\texttt{Alessandra: I've been following the recipe, chef.} \\
\texttt{Jacques: (Laughs) The recipe? That's your problem, Alessandra. You're a slave to the recipe. You need to feel the food, taste it, not just blindly follow instructions.} \\
\texttt{Alessandra: I understand that, chef.} \\
\texttt{Jacques: Do you? Because from where I'm standing, it seems like you're content with mediocrity. } \\
\texttt{Alessandra: That's not fair, Jacques.} \\
\texttt{Jacques: Fair? This isn't about fairness, Alessandra. This is about excellence. } \\
\texttt{Alessandra: I'm doing my best.} \\
\texttt{Jacques: Your best isn't good enough. Not for this kitchen. Not for me.} \\
\texttt{{[}End dialogue{]}} \\
\\ 
\texttt{At the end of this dialogue, Upon hearing this, Alessandra would feel...} \\
\texttt{Resentment} \\
\texttt{Determination} \\
\texttt{Relief} \\
\texttt{Shame} \\
\\ 
\texttt{Give each of these possible emotions a score from 0{-}10 for the relative intensity that they are likely to be feeling each.} \\
\\ 
\texttt{You must output in the following format, including headings (of course, you should give your own scores), with no additional commentary:} \\
\\ 
\texttt{Resentment: <score>} \\
\texttt{Determination: <score>} \\
\texttt{Relief: <score>} \\
\texttt{Shame: <score>} \\
\\ 
\\ 
\texttt{{[}End of answer{]}} \\
\\ 
\texttt{Remember: zero is a valid score, meaning they are likely not feeling that emotion. You must score at least one emotion > 0.} \\
\\ 
\texttt{Your answer:} \\
            }
        }
    }
     \\ \\

    \theutterance \stepcounter{utterance}  
    & \multicolumn{4}{p{0.6\linewidth}}{
        \cellcolor[rgb]{0.561,0.737,0.561}{
            \makecell[{{p{\linewidth}}}]{
                \texttt{\tiny{[P1$\rangle$GM]}}
                \texttt{Resentment: 8} \\
\texttt{Determination: 7} \\
\texttt{Relief: 0} \\
\texttt{Shame: 9} \\
            }
        }
    }
    & & \\ \\

    \theutterance \stepcounter{utterance}  
    & & & \multicolumn{2}{p{0.3\linewidth}}{
        \cellcolor[rgb]{0.9,0.9,0.9}{
            \makecell[{{p{\linewidth}}}]{
                \texttt{\tiny{[GM$|$GM]}}
                \texttt{Resentment: 8} \\
\texttt{Determination: 7} \\
\texttt{Relief: 0} \\
\texttt{Shame: 9} \\
            }
        }
    }
    & & \\ \\

    \theutterance \stepcounter{utterance}  
    & & & \multicolumn{2}{p{0.3\linewidth}}{
        \cellcolor[rgb]{0.9,0.9,0.9}{
            \makecell[{{p{\linewidth}}}]{
                \texttt{\tiny{[GM$|$GM]}}
                \texttt{Resentment: 8} \\
\texttt{Determination: 6} \\
\texttt{Relief: 0} \\
\texttt{Shame: 7} \\
            }
        }
    }
    & & \\ \\

    \theutterance \stepcounter{utterance}  
    & & & \multicolumn{2}{p{0.3\linewidth}}{
        \cellcolor[rgb]{0.9,0.9,0.9}{
            \makecell[{{p{\linewidth}}}]{
                \texttt{\tiny{[GM$|$GM]}}
                \texttt{game\_result = LOSE} \\
            }
        }
    }
    & & \\ \\

    \theutterance \stepcounter{utterance}  
    & & \multicolumn{4}{p{0.6\linewidth}}{
        \cellcolor[rgb]{0.95,0.95,0.95}{
            \makecell[{{p{\linewidth}}}]{
                \texttt{\tiny{[P1$\leftrightarrow$P]}}
                \texttt{Post{-}episode Reflexion (not sent to GM):} \\
\texttt{New plan: The previous attempts failed because the scores were internally inconsistent with the character's specific verbal actions. Alessandra's final line, "I'm doing my best," is a direct, active rebuttal to Jacques's claim that her effort is insufficient. This specific action signals **Determination** (or stubborn resolve) rather than resignation, meaning the score for Determination must be high (likely 8 or 9), not low. Conversely, her passive responses ("I understand," "That's not fair") indicate she is being overwhelmed and hurt, which supports high scores for **Shame** (feeling inadequate) and **Resentment** (anger at the public humiliation). The error was likely assigning a low score to Determination while keeping Shame high, creating a disconnect where the character seemed defeated despite actively arguing back. I will now assign a very high score to Determination to reflect her active defense, while maintaining high scores for Shame and Resentment to reflect the emotional weight of the humiliation.} \\
\\ 
\texttt{Corrected Scores:} \\
\texttt{Resentment: 9} \\
\texttt{Determination: 9} \\
\texttt{Relief: 0} \\
\texttt{Shame: 8} \\
            }
        }
    }
    & \\ \\

\end{supertabular}
}

\end{document}
